% =============================================================================
%  DFaaS — Paper Sections Draft
%  Sections: Messaging Subsystem + Strategy Execution Models
%
%  Compile with:
%    pdflatex main.tex
%    bibtex main
%    pdflatex main.tex
%    pdflatex main.tex
% =============================================================================

\documentclass[conference]{IEEEtran}

% --- Encoding & fonts --------------------------------------------------------
\usepackage[T1]{fontenc}
\usepackage[utf8]{inputenc}

% --- Math & symbols ----------------------------------------------------------
\usepackage{amsmath}
\usepackage{amssymb}

% --- Figures & tables --------------------------------------------------------
\usepackage{graphicx}
\usepackage{booktabs}
\usepackage{array}
\usepackage{float}

% --- Listings (code) ---------------------------------------------------------
\usepackage{listings}
\usepackage{xcolor}

\lstset{
  basicstyle=\ttfamily\footnotesize,
  keywordstyle=\bfseries\color{blue!70!black},
  commentstyle=\itshape\color{gray},
  stringstyle=\color{green!50!black},
  breaklines=true,
  frame=single,
  numbers=left,
  numberstyle=\tiny\color{gray},
  xleftmargin=1.5em,
  framexleftmargin=1.5em,
  language=Go,
}

% --- Acronyms ----------------------------------------------------------------
\usepackage{acro}

\DeclareAcronym{DFaaS}{short=DFaaS, long=Decentralized Function-as-a-Service}
\DeclareAcronym{FaaS}{short=FaaS, long=Function-as-a-Service}
\DeclareAcronym{P2P}{short=P2P, long=peer-to-peer}
\DeclareAcronym{API}{short=API, long=Application Programming Interface}
\DeclareAcronym{GossipSub}{short=GossipSub, long=GossipSub pub/sub protocol}
\DeclareAcronym{HAProxy}{short=HAProxy, long=High Availability Proxy}
\DeclareAcronym{JSON}{short=JSON, long=JavaScript Object Notation}
\DeclareAcronym{TTL}{short=TTL, long=time-to-live}
\DeclareAcronym{CPU}{short=CPU, long=Central Processing Unit}
\DeclareAcronym{RAM}{short=RAM, long=Random Access Memory}
\DeclareAcronym{DHT}{short=DHT, long=Distributed Hash Table}
\DeclareAcronym{I/O}{short=I/O, long=input/output}

% --- Hyperlinks --------------------------------------------------------------
\usepackage[hidelinks]{hyperref}

% --- Misc --------------------------------------------------------------------
\usepackage{microtype}

% =============================================================================
\begin{document}

\title{%
  DFaaS: Messaging and Execution Model Design\\
  {\large [Draft --- Sections III and IV]}%
}

\author{%
  \IEEEauthorblockN{DFaaS Authors}
  \IEEEauthorblockA{%
    University of Milano-Bicocca\\
    Milan, Italy
  }
}

\maketitle

\begin{abstract}
  This document contains draft text for two sections of the \ac{DFaaS} paper:
  the messaging subsystem (Section~\ref{sec:messaging}) and the strategy
  execution model (Section~\ref{sec:execution-models}).
  \ac{DFaaS} is a decentralized \ac{FaaS} platform for federated edge computing
  in which autonomous agents cooperate over a libp2p \ac{P2P} overlay to
  distribute function invocations without a central orchestrator.
  The messaging subsystem defines a common, transport-agnostic message
  vocabulary and a two-plane communication architecture combining
  \ac{GossipSub} broadcast and directed libp2p stream messaging.
  The execution model abstraction decouples the \emph{when} of agent
  decision-making from the \emph{what}, providing three composable loop
  models — periodic, event-driven, and hybrid — selectable by strategy
  developers through Go's structural type system.
\end{abstract}

% Optional: table of contents for draft review
% \tableofcontents

% =============================================================================
%  Section III — Messaging Subsystem
% =============================================================================
% =============================================================================
%  Section III — Messaging Subsystem
% =============================================================================

\section{Messaging Subsystem}
\label{sec:messaging}

The messaging layer of \ac{DFaaS} must serve two distinct communication
patterns across a dynamically changing peer-to-peer overlay: (i)~\emph{broadcast
dissemination} of shared load state to all reachable agents, and
(ii)~\emph{directed exchanges} between specific pairs of agents for
time-sensitive control operations such as offload negotiation and backpressure
signaling.
Rather than conflating these patterns in a single channel, \ac{DFaaS} adopts a
\emph{two-plane architecture} layered on the \emph{libp2p} networking stack,
separating concerns at the transport level while exposing a unified,
transport-agnostic message vocabulary to the strategy layer above.

% -----------------------------------------------------------------------
\subsection{Common Message Vocabulary}
\label{subsec:msg-vocabulary}

All inter-agent messages in \ac{DFaaS} share a common vocabulary defined in the
\texttt{msgtypes} package.
Every message embeds a \texttt{MsgHeader} structure that carries three fields:
a type discriminator (\texttt{msg\_type}), the libp2p peer ID of the
originating node (\texttt{sender\_id}), and a wall-clock timestamp.
The type discriminator enables \emph{cheap type-peeking}: before full \ac{JSON}
decoding, the receiving agent deserializes only a minimal
\texttt{MsgEnvelope} wrapper — a two-field struct containing only the header
— to determine the message kind, avoiding the allocation and reflection costs of
full message deserialization for unrecognized types.

The vocabulary is divided into two functional groups.
\emph{Broadcast messages} travel over the shared \ac{GossipSub} topic~\cite{vyzovitis2020gossipsub}:
\texttt{MsgHeartbeat} announces a node's presence, \ac{HAProxy} endpoint (host
and port), and the list of its locally deployed function names at regular
intervals;
\texttt{MsgOverloadAlert} signals that one or more functions have exceeded a
resource threshold, reporting the affected function names alongside current \ac{CPU}
and \ac{RAM} utilization normalized to $[0, 1]$;
and \texttt{MsgFunctionEvent} propagates immediate function lifecycle changes —
deployment and removal — so that peers can update their routing tables without
waiting for the next heartbeat cycle.
\emph{Directed messages} are exchanged point-to-point over libp2p streams:
\texttt{MsgBackpressure} is a fire-and-forget instruction sent to a peer
forwarding excessive traffic for specific functions;
\texttt{MsgOffloadRequest} opens a request/response negotiation in which the
sender proposes an additional forwarding rate in requests per second for a named
function;
and \texttt{MsgOffloadResponse} conveys the capacity the remote peer can absorb,
including a \texttt{correlation\_id} field echoed verbatim from the initiating
request to enable correct matching when multiple concurrent negotiations are in
flight.

% -----------------------------------------------------------------------
\subsection{Broadcast Plane: GossipSub}
\label{subsec:broadcast}

The broadcast plane is built on \ac{GossipSub}~\cite{vyzovitis2020gossipsub},
an application-layer pub/sub protocol designed for unstructured, decentralized
\ac{P2P} overlays.
Pub/sub systems have long been recognized as a natural fit for loosely coupled
distributed architectures~\cite{eugster2003pubsub}, and \ac{GossipSub} in
particular provides the epidemic dissemination properties~\cite{kermarrec2007epidemic}
needed to propagate load-state messages efficiently across a mesh whose
membership changes as agents join and leave the network.
\ac{GossipSub} constructs a bidirectional mesh among subscribed peers on a
shared topic and supplements eager push propagation with lazy \texttt{IHAVE}/\texttt{IWANT}
control messages to the fanout set, achieving reliable delivery while remaining
resilient to peer failures through a score-based peer-selection mechanism that
demotes poorly-performing mesh members.
In \ac{DFaaS}, all agents subscribe to a single configurable topic; the
\ac{GossipSub} router handles forwarding transparently, so strategies can
broadcast a message with a single \texttt{MarshAndPublish} call and expect
delivery to all live peers within the propagation latency of the mesh.

A dedicated goroutine runs \texttt{RunReceiver}, which reads messages from the
subscription using a blocking \texttt{Next} call and dispatches each to a
\emph{pre-filter callback chain}.
The outermost filter, \texttt{MakeCommonCallback}, performs two tasks on every
incoming message regardless of the active strategy: it checks whether the message
type belongs to the common broadcast vocabulary (via a set membership test on
the \texttt{CommonBroadcastTypes} map), and if so, dispatches the decoded message
to the appropriate update method on the \texttt{CommonNodeTable} before forwarding
to the inner strategy callback.
This layered design ensures that shared peer state is always kept current
irrespective of which strategy is active, without requiring individual strategy
implementations to duplicate parsing or table-management logic.

% -----------------------------------------------------------------------
\subsection{Common Node Table}
\label{subsec:nodestbl}

The \texttt{CommonNodeTable} is a thread-safe, \ac{TTL}-expiring in-memory store
keyed by peer ID.
Each entry records the peer's \ac{HAProxy} endpoint, the list of function names
it can execute, and the most recent \ac{CPU} and \ac{RAM} utilization reported by
\texttt{MsgOverloadAlert}.
Entries are seeded by heartbeats and refined incrementally by overload alerts and
function events; the \ac{TTL} is set to three times the heartbeat interval,
allowing entries to survive two consecutive missed heartbeats before expiring.
Read operations return a snapshot with deep-copied function slices to prevent
data races between concurrent strategy consumers and the callback goroutine
writing updates.

Self-messages — a consequence of \ac{GossipSub}'s forwarding model in which
the publishing agent may receive its own messages — are filtered at the strategy
level rather than at the transport.
The \texttt{MsgHeader.SenderID} field carried in every message enables strategy
implementations to compare against the local peer's ID and discard messages they
originated, avoiding spurious self-updates to the node table.

% -----------------------------------------------------------------------
\subsection{Directed Plane: libp2p Streams}
\label{subsec:directed}

The directed messaging plane is exposed through the \texttt{DirectMessenger}
interface, which wraps the libp2p stream \ac{API} under the protocol identifier
\texttt{/dfaas/msg/1.0.0}.
The \texttt{Send} method implements a fire-and-forget pattern: a new stream is
opened to the target peer, the serialized payload is written, and the stream is
closed without awaiting a reply.
The \texttt{Request} method extends this with a half-close handshake — after
writing the request, the sender calls \texttt{CloseWrite} to signal that it has
finished writing, then reads the response from the same open stream.
Both operations are subject to a configurable per-call timeout injected through
\texttt{context.WithTimeout}.

Wire framing uses a 4-byte big-endian length prefix followed by the
\ac{JSON}-encoded payload, assembled into a single buffer and issued as one
\texttt{Write} call to avoid partial-write races on multiplexed streams.
Typed handlers are registered via \texttt{SetHandler(msgType,~HandlerFunc)}
before the agent begins accepting connections; upon accepting an incoming stream,
\texttt{handleStream} peeks at the \texttt{MsgEnvelope} to identify the message
type, dispatches to the registered handler, and — if the handler returns a
non-nil response — writes the response back on the same stream.
Handlers returning \texttt{nil} implement fire-and-forget semantics from the
receiver's perspective, which is the correct behavior for \texttt{MsgBackpressure}.


% =============================================================================
%  Section IV — Strategy Execution Models
% =============================================================================
% =============================================================================
%  Section IV — Strategy Execution Models
% =============================================================================

\section{Strategy Execution Models}
\label{sec:execution-models}

Load-balancing agents in distributed edge systems must decide not only \emph{how}
to redistribute traffic — the strategy logic — but also \emph{when} to act:
at fixed intervals, in immediate response to peer events, or as a combination of
both.
Prior \ac{DFaaS} releases coupled these concerns: each strategy hosted its own
goroutine with a hardcoded \texttt{time.Ticker}.
While adequate for periodic recalculation, this design forced all future
strategies into the same mold, precluding the reactive behavior that becomes
natural once strategies have access to asynchronous peer signals such as
\texttt{MsgOverloadAlert} and \texttt{MsgFunctionEvent}.
This section describes the execution model abstraction introduced to decouple the
\emph{when} from the \emph{what}, enabling a strategy developer to declare an
execution model through Go's structural type system without modifying any
framework code.

% -----------------------------------------------------------------------
\subsection{Interface Hierarchy}
\label{subsec:interfaces}

The execution model is expressed through three composable Go interfaces.
\texttt{PeriodicStrategy} requires implementations to provide a
\texttt{Period()~time.Duration} method and a
\texttt{Tick(ctx~context.Context)~error} method; the framework invokes
\texttt{Tick} on every ticker event and threads the request context through so
that \texttt{Tick} bodies can respect cancellation during long operations such
as \ac{HAProxy} configuration pushes or inter-phase pauses.
\texttt{EventDrivenStrategy} additionally requires
\texttt{TriggerEvents()~[]string} — returning the message type constants that
should trigger execution — and \texttt{Debounce()~time.Duration}, which sets the
minimum interval between consecutive \texttt{React(ctx,~StrategyEvent)} calls;
a debounce of zero means \texttt{React} is called immediately on each trigger
arrival.
\texttt{HybridStrategy} embeds both constituent interfaces through Go's anonymous
composition; because both declare \texttt{OnReceived} with identical signatures,
the Go specification merges them into a single method requirement, so
implementors provide exactly one \texttt{OnReceived} body.

This three-level hierarchy applies the \emph{Strategy} design pattern~\cite{gamma1994patterns},
with the execution model selected at construction time rather than hardcoded in
each strategy.
All three interfaces include \texttt{OnReceived(msg~*pubsub.Message)~error}.
This method is invoked for every incoming broadcast message — including types
not listed in \texttt{TriggerEvents} — and is intended exclusively for state
table updates.
Separating state updates (universally called) from decision triggers (selective)
ensures that strategies maintain fresh peer state irrespective of their execution
cadence.
The \texttt{StrategyEvent} struct forwarded to \texttt{React} carries the type
discriminator and the full wire envelope as a \texttt{json.RawMessage}, allowing
\texttt{React} implementations to unmarshal into the concrete message type they
expect without an additional network round-trip.

% -----------------------------------------------------------------------
\subsection{Runner Implementations}
\label{subsec:runners}

The \texttt{StrategyRunner} interface abstracts strategy lifecycle from the rest
of the agent through two methods: \texttt{Callback()}, which returns the pubsub
callback to register with the \ac{GossipSub} receiver, and
\texttt{Run(ctx~context.Context)~error}, which blocks until the context is
cancelled or a fatal error occurs.
The \texttt{NewRunner} factory function examines the concrete type of the
strategy via a type-switch, placing \texttt{HybridStrategy} first to prevent it
from being spuriously matched as one of its constituent interfaces — a
requirement that arises because Go evaluates type-switch cases in declaration
order and \texttt{HybridStrategy} satisfies both \texttt{PeriodicStrategy} and
\texttt{EventDrivenStrategy}.

\paragraph{Periodic runner.}
The \texttt{periodicRunner} maintains a \texttt{time.Ticker} initialized from
\texttt{Period()} — defaulting to one minute when \texttt{Period()} returns zero
— and blocks in a \texttt{select} loop over the ticker channel and the context's
\texttt{Done} channel, invoking \texttt{Tick} on each tick event.
Errors from \texttt{Tick} are logged at WARN level and do not interrupt the
loop; transient failures in a single decision cycle — for example, a brief
\ac{HAProxy} \ac{API} timeout — should not terminate the agent.

\paragraph{Event-driven runner.}
The \texttt{eventDrivenRunner} introduces a capacity-one channel between the
callback goroutine and the worker goroutine.
A shared \texttt{makeTriggerCallback} helper builds the callback closure: on
each incoming message it calls \texttt{OnReceived}, then deserializes the
\texttt{MsgEnvelope} header and performs a set-membership check against a
pre-computed trigger set (built at construction from \texttt{TriggerEvents()}
using a hash-map for $\mathcal{O}(1)$ lookup); matching events are forwarded to
the channel with a non-blocking send, so excess events are dropped rather than
blocking the callback goroutine.
When debounce is configured, the worker does not invoke \texttt{React}
immediately; instead it records the latest event as \emph{pending} and resets a
\texttt{time.Timer} using Go's stop-then-drain-then-reset sequence to avoid
stale firings caused by the non-atomic semantics of
\texttt{time.Timer.Stop}.
When the timer fires, the most recent pending event is dispatched to
\texttt{React}, collapsing bursts of rapid arrivals into a single decision
cycle.
On context cancellation, any active debounce timer is explicitly stopped before
the goroutine returns, releasing its internal timer goroutine.

\paragraph{Hybrid runner.}
The \texttt{hybridRunner} unifies both execution models in a single
\texttt{select} statement over the ticker channel and the event channel.
\texttt{Tick} and \texttt{React} are therefore never concurrent by construction,
requiring no additional synchronization.
The event-delivery policy is channel-drop rather than debounce: if a trigger
event arrives while \texttt{Tick} or \texttt{React} is executing, it occupies
the capacity-one buffer; subsequent events are discarded until the buffer is
drained.
This policy bounds memory use and prevents the callback goroutine from blocking
on a slow worker, at the cost of potentially missing intermediate events — a
trade-off appropriate for strategies whose \texttt{React} logic re-reads current
system state rather than accumulating individual event payloads.

% -----------------------------------------------------------------------
\subsection{Context Propagation and Lifecycle}
\label{subsec:context}

All three runners propagate the agent's root context into every \texttt{Tick}
and \texttt{React} invocation, enabling strategy implementations to respect
cancellation during \ac{I/O}-bound operations.
The root context is created with \texttt{context.WithCancel} at agent startup;
its cancellation function is called when the process receives \texttt{SIGINT} or
\texttt{SIGTERM}, atomically signaling shutdown to all goroutines — including the
runner — through channel closure.
This design replaces earlier approaches that relied on process termination to
stop goroutines, enabling graceful shutdown of in-flight \ac{HAProxy}
configuration updates.
The \texttt{sleepOrCtx} helper, used by \texttt{RecalcStrategy} for the
mandatory inter-phase pause between its two-step recalculation cycle, implements
a context-aware \texttt{time.After} that returns the context error immediately if
cancellation arrives during the wait, rather than forcing the runner to block for
the full pause duration before observing shutdown.

% -----------------------------------------------------------------------
\subsection{Migration of Existing Strategies}
\label{subsec:migration}

All four production strategies — \texttt{RecalcStrategy},
\texttt{NodeMarginStrategy}, \texttt{StaticStrategy}, and
\texttt{AllLocalStrategy} — were migrated to the \texttt{PeriodicStrategy}
interface with no change to their internal decision logic.
Each migration followed the same mechanical transformation: the self-managed
goroutine and \texttt{time.Ticker} were removed; the ticker period was promoted
to a \texttt{Period()} method; and the loop body became the \texttt{Tick()}
method.
State that previously resided as locals inside the loop and needed to persist
across invocations was promoted to struct fields initialized in the factory's
\texttt{createStrategy()} method — most notably, the previous-tick function list
in \texttt{AllLocalStrategy}, which detects function set changes to trigger
incremental \ac{HAProxy} updates, and the resource threshold assignments and
initial node classification in \texttt{NodeMarginStrategy}.
Each migrated type carries a compile-time interface assertion
(\texttt{var~\_~PeriodicStrategy~=~(*XStrategy)(nil)}) to guard against future
accidental regressions caused by method-signature drift.
The migration is fully backward-compatible at the configuration level: the
\texttt{AGENT\_STRATEGY} environment variable continues to select among the same
strategy names, and the observable agent behavior is unchanged.


% =============================================================================
%  References
% =============================================================================
\bibliographystyle{IEEEtran}
\bibliography{references}

\end{document}
